Investors and regulators face difficulty distinguishing implemented AI capability from narrative positioning in annual 10-K filings. Using U.S. public firms from 2016 to 2025, we classify AI-related sentences in corporate filings as actionable, speculative, or irrelevant and aggregate them into filing-based disclosure measures. We combine low-credibility disclosure composition with contemporaneous AI patent weakness to construct PatentMismatch, a scalable and audited proxy for low-credibility AI disclosure. After 2022, PatentMismatch characterizes more than 40\% of AI-talking firm-years in both 2024 and 2025. PatentMismatch predicts weaker subsequent AI patent grants and profitability, and higher subsequent R\&D intensity. Disclosure credibility appears stronger following SEC comment-letter scrutiny, weaker around large equity-issuance windows, and lower when CEO equity incentives are stronger. Filing-date reactions remain muted. The evidence points to a disclosure-credibility problem whose strongest patterns appear in later realization, scrutiny, incentives, and financing-window dynamics rather than in immediate market pricing.